\chapter{PENELITIAN TERKAIT}

\section{Efisiensi Inferensi dan Komputasi Adaptif}

Model Transformer memproses masukan melalui sejumlah lapisan, dan pada model decoder-only proses tersebut diulang untuk setiap token keluaran. Akibatnya, inferensi dapat menjadi mahal, terutama pada perangkat dengan sumber daya terbatas. Berbagai penelitian efisiensi inferensi berupaya mengurangi komputasi ini sambil mempertahankan kemampuan model menyelesaikan tugasnya.

Salah satu cara mengurangi komputasi adalah menjalankan jaringan dengan kedalaman yang lebih pendek. LayerDrop menggunakan \textit{structured dropout} selama pelatihan agar subjaringan dengan kedalaman berbeda dapat digunakan saat inferensi dengan penurunan kinerja yang terbatas \citep{fan2020layerdrop}. Pendekatan lain, yaitu \textit{early exit}, menghentikan pemrosesan pada lapisan antara ketika keluaran dinilai telah cukup baik. Sebagai contoh, DeeBERT menambahkan titik keluar pada beberapa lapisan BERT sehingga sebagian masukan tidak perlu melewati seluruh jaringan \citep{xin2020deebert}.

Pada pembangkitan autoregresif, keputusan komputasi dapat dibuat pada tingkat urutan, token, atau lapisan. Confident Adaptive Language Modeling (CALM) mengalokasikan komputasi secara adaptif untuk setiap masukan dan langkah pembangkitan token. Pendekatan tersebut menunjukkan bahwa keputusan pada tingkat token memerlukan penanganan khusus terhadap dependensi antartoken dan representasi yang tidak dihitung pada langkah sebelumnya \citep{schuster2022calm}. LayerSkip juga mengkaji efisiensi model bahasa autoregresif melalui pelatihan dengan \textit{layer dropout} dan \textit{early-exit loss}, kemudian menggunakan keluaran lapisan awal untuk \textit{self-speculative decoding} \citep{elhoushi2024layerskip}.

Secara keseluruhan, pendekatan-pendekatan tersebut mengurangi komputasi dengan cara dan pada tahap yang berbeda; sebagian juga memerlukan pelatihan khusus pada model utama. Penelitian ini mengambil arah yang lebih terfokus: bobot dasar Qwen3-0.6B dibekukan ketika router dilatih, lalu adapter LoRA dilatih secara terpisah untuk mengompensasi dampak pelompatan lapisan. Perbedaan cara menentukan lapisan yang dilewati menjadi dasar pembahasan berikutnya mengenai strategi statis dan routing adaptif.

\section{Pelompatan Lapisan Statis dan Routing Adaptif}

Pelompatan lapisan dapat diterapkan dengan pola tetap atau dengan keputusan yang menyesuaikan masukan. Strategi statis menjalankan pola lapisan yang sama untuk semua artikel sehingga sederhana dan mudah direproduksi, tetapi tidak menyesuaikan komputasi dengan isi artikel. Sebaliknya, routing adaptif memilih jalur eksekusi berdasarkan karakteristik setiap masukan. Perbedaan ini relevan karena artikel dengan tingkat kesulitan dan struktur informasi yang berbeda mungkin memerlukan jumlah komputasi yang berbeda pula.

PuDDing menggunakan prompt untuk memilih pola kedalaman dari beberapa pola lapisan yang telah disiapkan. Hasilnya menunjukkan potensi routing berbasis prompt untuk mengurangi komputasi, tetapi pendekatan tersebut memilih di antara pola yang telah ditentukan dan dievaluasi pada tugas penalaran akal sehat \citep{wee2025pudding}. Penelitian ini menggunakan keluaran biner terpisah untuk setiap kelompok lapisan tengah. Dengan demikian, pola route tidak diperlakukan sebagai kelas dengan identitas tertentu dan jumlah lapisan aktif dapat berbeda antarartikel.

FiRST memperluas gagasan routing berbasis prompt dengan router khusus pada lapisan dan keputusan yang dibuat pada tahap \textit{prefill}, kemudian dipakai selama proses pembentukan token keluaran. FiRST melatih router lebih dahulu dengan backbone beku, lalu membekukan router dan melatih adapter LoRA pada modul attention dan FFN untuk memulihkan kualitas setelah pelompatan lapisan \citep{jain2025first}. Penelitian ini mengadopsi pola pelatihan dua tahap tersebut, tetapi tetap menggunakan satu router pada tingkat kelompok lapisan, bukan router terpisah untuk setiap lapisan. Target router penelitian ini dibentuk melalui pencarian bertahap berdasarkan kerugian prediksi token dan proporsi kelompok aktif.

Beberapa penelitian juga menunjukkan bahwa pengurangan komputasi pada model autoregresif harus memperhatikan cara representasi dan \textit{key-value cache} dikelola. FFN-SkipLLM melaporkan bahwa strategi \textit{early exit} yang tidak tepat dapat berkaitan dengan keruntuhan generasi, halusinasi, dan penurunan kualitas, lalu mengusulkan pelompatan blok feed-forward yang tidak memerlukan penanganan KV cache dengan cara yang sama \citep{jaiswal2024ffnskip}. Oleh karena itu, penelitian ini memasukkan pemeriksaan bahwa lapisan yang dilewati benar-benar tidak dieksekusi, route tidak berubah selama decoding, dan KV cache konsisten. Efisiensi utama juga diukur menggunakan waktu inferensi aktual, bukan hanya jumlah lapisan yang dilewati.

\section{Peringkasan Berita Bahasa Indonesia dan XL-Sum}

Untuk menguji perbedaan strategi tersebut pada tugas yang terukur, penelitian ini menggunakan XL-Sum, dataset peringkasan abstraktif multibahasa yang menyediakan pasangan artikel berita dan ringkasan acuan dalam 44 bahasa, termasuk bahasa Indonesia \citep{hasan2021xlsum}. Pasangan ini dapat digunakan baik untuk membentuk target pelatihan router maupun untuk menilai ringkasan hasil generasi. Penelitian dibatasi pada subset bahasa Indonesia, dengan prosedur pemrosesan data yang sama bagi seluruh metode pembanding.

Pada dataset tersebut, penelitian menggunakan Qwen3-0.6B, model bahasa kausal dengan 28 lapisan \citep{qwen2025modelcard}. Bobot dasarnya dibekukan selama pelatihan router, kemudian adapter LoRA dilatih pada tahap berikutnya. Rancangan evaluasi membandingkan model penuh, pola statis, dan routing adaptif dengan LoRA, serta menyertakan routing adaptif tanpa LoRA sebagai kontrol. Dengan demikian, perbandingan tidak hanya menguji strategi routing, tetapi juga memungkinkan pengamatan atas kontribusi LoRA.

\section{Evaluasi Otomatis Peringkasan}

ROUGE merupakan keluarga metrik yang membandingkan tumpang tindih unit teks antara ringkasan keluaran dan ringkasan acuan. ROUGE-N mengukur kecocokan n-gram, sedangkan ROUGE-L menggunakan subsekuens umum terpanjang \citep{lin2004rouge}. Penelitian ini menggunakan ROUGE-1, ROUGE-2, dan ROUGE-L untuk menangkap kecocokan kata, pasangan kata berurutan, dan struktur urutan teks.

BERTScore melengkapi ROUGE dengan membandingkan representasi kontekstual token menggunakan kemiripan kosinus \citep{zhang2020bertscore}. Metrik ini dapat menangkap kemiripan makna yang tidak selalu menghasilkan tumpang tindih kata secara langsung. Namun, kedua kelompok metrik tersebut tetap merupakan ukuran kecocokan terhadap acuan. Forde et al. menunjukkan bahwa metrik otomatis dapat kurang andal untuk menilai kualitas peringkasan multibahasa, dengan masalah yang lebih besar pada bahasa Indonesia \citep{forde2024multilingeval}. Karena itu, penelitian ini menggunakan beberapa metrik secara terpisah dan tidak menyimpulkan bahwa skor tersebut menjamin faktualitas atau kualitas ringkasan secara menyeluruh.

\section{Posisi Penelitian}

Uraian penelitian terdahulu memperlihatkan beragam cara mengurangi komputasi Transformer, mulai dari pelatihan agar model mampu menggunakan kedalaman lebih pendek, \textit{early exit}, hingga routing berbasis prompt. Meski memiliki tujuan serupa, pendekatan tersebut berbeda dalam tingkat pengambilan keputusan, kebutuhan perubahan pada model utama, tugas evaluasi, dan cara mengukur efisiensi. Perbedaan ini menunjukkan perlunya pengujian dalam konteks tugas dan rancangan yang spesifik.

Penelitian ini menguji \textit{sequence-level routing} pada Qwen3-0.6B untuk peringkasan berita bahasa Indonesia menggunakan XL-Sum. Router membaca representasi token artikel setelah kelompok pertama dan memilih secara biner apakah setiap kelompok dari lima kelompok lapisan tengah dijalankan atau dilewati. Keputusan yang sama dipertahankan selama sisa pemrosesan prompt dan seluruh proses pembentukan token keluaran. Router dilatih saat parameter Qwen dibekukan; setelah itu, router dibekukan dan LoRA dilatih sebagai tahap kompensasi. Routing adaptif tanpa LoRA menjadi kontrol untuk membandingkan hasil dengan konfigurasi yang menggunakan LoRA.

Evaluasi membandingkan routing adaptif dengan model penuh dan satu pola pelompatan statis yang dipilih menggunakan data validasi. LoRA dilatih pada masing-masing pendekatan utama, sedangkan routing adaptif tanpa LoRA menjadi kontrol. Perbandingan mencakup kualitas keluaran dan waktu inferensi aktual pada perangkat serta pengaturan yang sama, termasuk biaya router, \textit{prefill}, dan \textit{decoding}. Dengan demikian, penelitian menguji trade-off kualitas--efisiensi routing adaptif serta perubahan hasil ketika LoRA digunakan. Penelitian ini tidak mengklaim bahwa penggunaan router, pembagian kelompok lapisan, atau model Qwen3-0.6B tersebut merupakan penerapan pertama.
